\documentclass[conference]{IEEEtran}
\IEEEoverridecommandlockouts
\usepackage{cite}
\usepackage{amsmath,amssymb,amsfonts}
\usepackage{algorithmic}
\usepackage{graphicx}
\usepackage{textcomp}
\usepackage{xcolor}
\usepackage{hyperref}
\usepackage{orcidlink}
\usepackage[utf8]{inputenc}
\usepackage{float}
\def\BibTeX{{\rm B\kern-.05em{\sc i\kern-.025em b}\kern-.08em
    T\kern-.1667em\lower.7ex\hbox{E}\kern-.125emX}}
\begin{document}

\title{Protector analógico de sobrecorriente para motores DC\\}
\author{
    \IEEEauthorblockN{Ana Freitez\, \orcidlink{0009-0009-6506-499X}}
    \IEEEauthorblockA{\textit{Facultad de Ingeniería,} \\ 
    \textit{Licenciatura en Ingeniería Mecatrónica} \\ 
    \textit{Universidad Latina de Panamá} \\ 
    Ciudad de Panamá, Panamá}
    \and
    \IEEEauthorblockN{Cristell Cedeño\, \orcidlink{0009-0002-2617-663X}}
    \IEEEauthorblockA{\textit{Facultad de Ingeniería,} \\ 
    \textit{Licenciatura en Ingeniería Mecatrónica} \\ 
    \textit{Universidad Latina de Panamá} \\ 
    Ciudad de Panamá, Panamá}
}

\maketitle

\begin{abstract}
% Resumen breve (150-250 palabras): objetivo del proyecto, circuito implementado
% (integrador y derivador en cascada con amplificadores operacionales),
% metodología utilizada (simulación y/o implementación física) y principales
% resultados obtenidos.
[Completar resumen del proyecto aquí.]
\end{abstract}

\begin{IEEEkeywords}
[Completar}
\end{IEEEkeywords}

\section{Introducción}
% Contexto general del tema, motivación del proyecto y objetivo principal.
% Mencionar brevemente qué problema resuelve un circuito integrador/derivador
% en cascada y por qué es relevante estudiarlo.
[Completar introducción.]

\subsection{Objetivos}
\begin{itemize}
    \item Objetivo general: [Completar]
    \item Objetivos específicos:
    \begin{itemize}
        \item Completar
        \item Completar
    \end{itemize}
\end{itemize}

\section{Marco Teórico}

[Completar.]

\section{Metodología}

\subsection{Diseño del Circuito}
[Completar descripción del diseño.]

\subsection{Materiales y Equipo}
\begin{table}[htbp]
\caption{Lista de Materiales y Equipos}
\label{tab:materiales}
\centering
\begin{tabular}{l l c}
\hline
\textbf{Componente} & \textbf{Especificación} & \textbf{Cantidad} \\
\hline
[Completar.] & [Completar.] & [Completar.]\\
[Completar.] & [Completar.] & [Completar.]\\
[Completar.] & [Completar.] & [Completar.]\\
\hline
\end{tabular}
\end{table}

\subsection{Procedimiento}
[Completar.] & [Completar.] & [Completar.]\\

\section{Resultados}

\subsection{Simulación}

[Completar descripción de los resultados de simulación.]

\begin{figure}[H]
	\centering
	\includegraphics[width=0.7\linewidth]{Esquematico.png}
	\caption{Esquemático simulado en Proteus 8 Professional}
	\label{fig:esquematico}
\end{figure}
\begin{figure}[H]
	\centering
	\includegraphics[width=0.7\linewidth]{Simulacion_osciloscopio.png}
	\caption{Formas de onda obtenidas en el osciloscopio virtual: señal de entrada, salida del integrador y salida del derivador.}
	\label{fig:sim_osc}
\end{figure}

\subsection{Implementación Física}

\begin{figure}[H]
	\centering
	\includegraphics[width=0.7\linewidth]{Montaje_fisico.jpeg}
	\caption{Montaje físico del circuito en protoboard}
	\label{fig:montaje_fisico}
\end{figure}

\begin{figure}[H]
	\centering
	\includegraphics[width=0.7\linewidth]{Captura_osciloscopio_real.png}
	\caption{Funcionamiento real de la protección}
	\label{fig:capturaosciloscopioreal}
\end{figure}

\subsection{Datos Medidos}
% Completar la tabla con los valores obtenidos en simulación y/o en el
% laboratorio para distintas frecuencias de entrada.
\begin{table}[htbp]
\caption{Amplitudes de Salida Medidas en Función de la Frecuencia de Entrada}
\label{tab:datos}
\centering
\begin{tabular}{c c c c}
\hline
\textbf{$f$ (Hz)} & \textbf{$V_{in}$ (V)} &
\textbf{$V_{out,int}$ (V)} & \textbf{$V_{out,der}$ (V)}\\
\hline
25    & 4 & 4.400 V & 444.0 mV \\
50    & 4 & [Completar] & [Completar]\\
100   & 4 & [Completar] & [Completar] \\
200   & 4 & [Completar] & [Completar] \\
500   & 4 & [Completar] & [Completar] \\
1000  & 4 & [Completar] & [Completar] \\
\hline
\end{tabular}
\end{table}

\section{Análisis de Resultados}
% Comparar los resultados de simulación con los resultados físicos.
% Explicar cualquier diferencia observada (ruido, saturación, tolerancias
% de componentes, etc.) y relacionarlos con la teoría de la Sección II.
[Completar análisis comparando teoría, simulación e implementación física.]

\section{Conclusiones}
% Resumir los hallazgos principales y si se cumplieron los objetivos
% planteados en la introducción.
[Completar conclusiones.]

\begin{thebibliography}{00}
\bibitem{b1} [Completar.]
\bibitem{b2} [Completar.]
\bibitem{b3} [Completar.]
\end{thebibliography}

\end{document}
